\documentclass[12pt,a4paper]{article}

\usepackage[T1]{fontenc}
\usepackage[utf8]{inputenc}
\usepackage[french]{babel}
\usepackage{lmodern}
\usepackage{geometry}
\usepackage{xcolor}
\usepackage{booktabs}
\usepackage{array}
\usepackage{hyperref}
\usepackage{listings}
\usepackage{graphicx}
\usepackage{float}

\geometry{margin=2.5cm}
\hypersetup{colorlinks=true,linkcolor=blue,urlcolor=blue}

\definecolor{codegray}{RGB}{245,245,245}
\definecolor{codegreen}{RGB}{35,110,55}
\definecolor{codeblue}{RGB}{25,70,145}

\lstset{
  backgroundcolor=\color{codegray},
  basicstyle=\ttfamily\small,
  keywordstyle=\color{codeblue}\bfseries,
  commentstyle=\color{codegreen},
  stringstyle=\color{red!65!black},
  breaklines=true,
  frame=single,
  showstringspaces=false,
  columns=fullflexible
}

\title{\textbf{TP 2 -- Intégration de la sécurité\ dans un pipeline CI/CD}}
\author{Étudiant : \rule{7cm}{0.4pt}}
\date{19 octobre 2026}

\begin{document}

\maketitle
\tableofcontents
\newpage

\section{Introduction}

Ce travail pratique a pour objectif d'intégrer des contrôles de sécurité automatisés
dans un pipeline d'intégration continue. L'idée principale est de transformer le
pipeline CI/CD en garde-fou : une modification ne doit pas être considérée comme
acceptable si elle introduit une vulnérabilité détectable par les outils de sécurité.

Le projet utilisé est une petite application Flask exposée sur le port 3000. Le dépôt
GitHub utilisé est :
\begin{center}
\url{https://github.com/ayabelkacem/Tp-2}
\end{center}

Les contrôles intégrés sont les suivants :
\begin{itemize}
  \item tests unitaires avec pytest ;
  \item SAST avec Bandit ;
  \item SCA avec Trivy ;
  \item DAST avec OWASP ZAP dans un conteneur Docker temporaire.
\end{itemize}

\section{Application et structure du projet}

L'application retourne le texte \texttt{Hello} sur la route principale. Une route
\texttt{/ping} est également présente dans la version corrigée. Elle valide le nom
d'hôte fourni par l'utilisateur avant d'appeler la commande système \texttt{ping}.

Les principaux fichiers du projet sont :
\begin{table}[H]
\centering
\begin{tabular}{p{4cm}p{9cm}}
\toprule
\textbf{Fichier} & \textbf{Rôle} \\
\midrule
\texttt{app.py} & Application Flask et routes HTTP. \\
\texttt{test\_app.py} & Test de la route principale. \\
\texttt{requirements.txt} & Versions de Flask et pytest. \\
\texttt{Dockerfile} & Image Docker de l'application. \\
\texttt{.github/workflows/ci.yml} & Pipeline GitHub Actions. \\
\texttt{.zap/rules.tsv} & Règles de sévérité utilisées par OWASP ZAP. \\
\bottomrule
\end{tabular}
\caption{Structure fonctionnelle du projet}
\end{table}

Le fichier \texttt{requirements.txt} contient les dépendances principales :
\begin{lstlisting}[language={},caption={Dépendances Python}]
Flask==3.0.3
pytest==8.3.3
\end{lstlisting}

\section{Pipeline CI/CD de base}

Le pipeline est déclenché dans deux cas : lors d'un \textit{push} sur \texttt{main}
et lors de la création ou de la mise à jour d'une pull request vers \texttt{main}.

Le premier job, nommé \texttt{build-test}, utilise Ubuntu et Python 3.12. Il installe
les dépendances puis exécute pytest :
\begin{lstlisting}[language={},caption={Job de compilation logique et de tests}]
build-test:
  runs-on: ubuntu-latest
  steps:
    - uses: actions/checkout@v4
    - uses: actions/setup-python@v5
      with:
        python-version: "3.12"
    - run: pip install -r requirements.txt
    - run: pytest
\end{lstlisting}

Le test réalisé vérifie que la route \texttt{/} retourne le statut HTTP 200. Ce job
constitue la première barrière du pipeline : les analyses de sécurité ne démarrent
qu'après la réussite des tests.

\section{Intégration du SAST avec Bandit}

Le SAST analyse le code source sans exécuter l'application. Bandit est installé dans
le job \texttt{sast}, puis lancé avec :
\begin{lstlisting}[language=bash]
bandit -r . -lll
\end{lstlisting}

L'option \texttt{-lll} demande à Bandit de considérer les problèmes de niveau élevé
comme bloquants. Le job dépend de \texttt{build-test} grâce à la directive :
\begin{lstlisting}[language={}]
needs: build-test
\end{lstlisting}

\subsection{Démonstration de l'échec}

Une branche \texttt{feat/vuln-demo} a été créée avec une route volontairement
vulnérable :
\begin{lstlisting}[language=Python]
import subprocess
from flask import request

@app.route("/ping")
def ping():
    host = request.args.get("host")
    return subprocess.check_output("ping -c 1 " + host, shell=True)
\end{lstlisting}

Cette implémentation concatène directement une entrée utilisateur dans une commande
exécutée par un shell. Elle permet une injection de commandes système. Bandit détecte
le problème \texttt{B602} et signale une sévérité \textbf{High} à cause de
\texttt{shell=True}.

La pull request correspondante est la PR \#1 :
\begin{center}
\url{https://github.com/ayabelkacem/Tp-2/pull/1}
\end{center}

Résultat observé : \texttt{build-test} réussit, mais \texttt{sast} échoue. Cette
branche démontre que le pipeline identifie et bloque une vulnérabilité de code.

\section{Intégration du SCA avec Trivy}

Le SCA analyse les dépendances et le système de fichiers à la recherche de
vulnérabilités connues. Le job Trivy utilise les paramètres suivants :
\begin{lstlisting}[language={}]
scan-type: fs
scanners: vuln
severity: HIGH,CRITICAL
ignore-unfixed: true
exit-code: "1"
\end{lstlisting}

La combinaison de \texttt{severity: HIGH,CRITICAL} et de \texttt{exit-code: "1"}
rend le job bloquant lorsqu'une vulnérabilité haute ou critique corrigible est
détectée. Les vulnérabilités sans correctif sont ignorées grâce à
\texttt{ignore-unfixed: true}.

\section{Analyse dynamique avec OWASP ZAP}

Le DAST teste l'application en fonctionnement. Le job \texttt{dast} dépend de la
réussite de SAST et de SCA :
\begin{lstlisting}[language={}]
needs: [sast, sca]
\end{lstlisting}

L'environnement est éphémère. Le pipeline construit l'image Docker, démarre un
conteneur nommé \texttt{app}, attend que l'URL soit disponible, puis lance ZAP :
\begin{lstlisting}[language=bash]
docker build -t app-under-test .
docker run -d --name app -p 3000:3000 app-under-test
curl -sf http://localhost:3000

docker run --rm --network host \\
  -v "$PWD:/zap/wrk:rw" \\
  ghcr.io/zaproxy/zaproxy:stable \\
  zap-baseline.py -t http://localhost:3000 \\
  -c .zap/rules.tsv -r zap-report.html -I
\end{lstlisting}

Le rapport HTML est téléversé comme artefact même en cas d'échec. Le conteneur est
également supprimé avec \texttt{if: always()} afin d'éviter de laisser un service
actif sur le runner.

\subsection{Règles ZAP}

Le fichier \texttt{.zap/rules.tsv} utilise de vraies tabulations entre les colonnes :
\begin{lstlisting}[language={}]
40012    FAIL    (Cross Site Scripting - Reflected)
40018    FAIL    (SQL Injection)
90020    FAIL    (Remote OS Command Injection)
10038    WARN    (Content Security Policy Header Not Set)
\end{lstlisting}

Les règles XSS, SQL Injection et Remote OS Command Injection sont bloquantes. La
règle CSP est un avertissement. L'option ZAP \texttt{-I} permet de ne pas bloquer le
pipeline pour les avertissements, tout en conservant le blocage des alertes
classées \texttt{FAIL}. Le rapport ZAP reste disponible comme artefact GitHub Actions.

\section{Correction de la vulnérabilité}

La branche \texttt{fix/vuln-demo} remplace l'exécution shell par une validation stricte
et une liste d'arguments :
\begin{lstlisting}[language=Python]
import re
import subprocess
from flask import request

@app.route("/ping")
def ping():
    host = request.args.get("host", "")
    if not re.fullmatch(r"[A-Za-z0-9.\-]{1,253}", host):
        return "Invalid host", 400
    result = subprocess.run(
        ["ping", "-c", "1", host],
        capture_output=True,
        text=True,
    )
    return result.stdout
\end{lstlisting}

La validation limite l'entrée à des caractères alphanumériques, aux points et au
tiret, avec une longueur maximale de 253 caractères. Surtout, la commande n'est plus
interprétée par un shell : l'option \texttt{shell=True} a disparu.

La pull request de correction est la PR \#2 :
\begin{center}
\url{https://github.com/ayabelkacem/Tp-2/pull/2}
\end{center}

Les résultats finaux de la PR corrigée sont :
\begin{itemize}
  \item \texttt{build-test} : succès ;
  \item \texttt{sast} : succès ;
  \item \texttt{sca} : succès ;
  \item \texttt{dast} : succès.
\end{itemize}

\section{Synthèse des résultats}

\begin{table}[H]
\centering
\begin{tabular}{p{3cm}p{3.1cm}p{3.1cm}p{4.2cm}}
\toprule
\textbf{Contrôle} & \textbf{PR vulnérable} & \textbf{PR corrigée} & \textbf{Interprétation} \\
\midrule
Build et tests & Succès & Succès & Le projet se construit et les tests passent. \\
SAST Bandit & Échec & Succès & L'injection shell est détectée puis supprimée. \\
SCA Trivy & Succès & Succès & Les dépendances ne présentent pas de blocage HIGH/CRITICAL. \\
DAST ZAP & Neutre & Succès & Le scan dynamique est exécuté après SAST et SCA. \\
\bottomrule
\end{tabular}
\caption{Comparaison des deux pull requests}
\end{table}

\section{Protection de la branche principale}

Pour que le pipeline soit réellement un garde-fou lors d'une fusion, les règles GitHub
suivantes doivent être activées pour la branche \texttt{main} :
\begin{enumerate}
  \item ouvrir \textit{Settings} puis \textit{Branches} ;
  \item créer une règle pour le motif \texttt{main} ;
  \item activer \textit{Require a pull request before merging} ;
  \item activer \textit{Require status checks to pass before merging} ;
  \item rendre obligatoires les contrôles \texttt{build-test}, \texttt{sast},
        \texttt{sca} et \texttt{dast}.
\end{enumerate}

Cette configuration est importante : sans elle, un utilisateur disposant des droits
appropriés pourrait fusionner une branche malgré un pipeline rouge.

\section{Captures d'écran}

Les captures suivantes peuvent être insérées dans cette section. Il faut remplacer les
noms de fichiers par les captures réellement exportées depuis GitHub :

\begin{figure}[H]
\centering
\IfFileExists{capture-pr-vulnerable.png}{%
  \includegraphics[width=0.95\textwidth]{capture-pr-vulnerable.png}%
}{\fbox{\parbox{0.85\textwidth}{\centering Capture à insérer : PR \#1 avec le job SAST en échec.}}}
\caption{Échec du pipeline sur la pull request vulnérable}
\end{figure}

\begin{figure}[H]
\centering
\IfFileExists{capture-pr-corrigee.png}{%
  \includegraphics[width=0.95\textwidth]{capture-pr-corrigee.png}%
}{\fbox{\parbox{0.85\textwidth}{\centering Capture à insérer : PR \#2 avec les quatre jobs au vert.}}}
\caption{Succès du pipeline sur la pull request corrigée}
\end{figure}

\section{Conclusion}

Le TP a permis de mettre en place une chaîne de contrôles de sécurité dans GitHub
Actions. Le pipeline vérifie d'abord le fonctionnement de l'application, puis analyse
le code avec Bandit, les dépendances avec Trivy et l'application en fonctionnement avec
OWASP ZAP.

La démonstration est complète : une vulnérabilité volontaire provoque l'échec du SAST,
puis la correction de cette vulnérabilité permet à la pull request de passer les quatre
contrôles. La dernière étape nécessaire pour imposer ce comportement à toute l'équipe
est l'activation effective de la protection de la branche \texttt{main} dans les
paramètres GitHub.

\end{document}